\chapter{Zenoh Pico wrapper}
\label{chap:zenoh-pico-wrapper}
\sidenote{non mi veniva un nome migliore}

Breve introduzione alla creazione di un wrapper safe-rust per zenoh pico ed alle
varie sezioni.

\section{L'ownership model in Zenoh Pico}
\label{sec:ownership-model-in-zenoh-pico}
\sidenote{Parlerei qui di questa cosa invece che nel capitolo di zenoh pico essendo
strettamente collegato a come è stato implementato il wrapper}

Come zenoh pico simuli l'ownership model di Rust in C per mantenere un API il più
simile possibile alla libreria originale e per facilitare la gestione della memoria.

\section{Il wrapper}
\label{sec:the-wrapper}

\subsection{Le macro in Rust}
\label{sec:macros-in-rust}

Introduzione alle macro in Rust ed alla generazione di codice a tempo di compilazione.

\subsection{L'implementazione tramite macro}
\label{sec:implementation-via-macros}

Come il wrapper sia stato implementato tramite macro sfruttando la strutturazione
simil-rust della libreria originale \cref{sec:ownership-model-in-zenoh-pico}
